\section{Resumen Ejecutivo}
\label{sec:resumen-ejecutivo}

Este informe documenta AeroTrack, un sistema de visión computacional que estima la
concentración de personas por zona y sigue su recorrido entre cámaras sobre el plano de un
espacio, pensado para el Aeropuerto Internacional Jorge Chávez a partir de sus cámaras fijas.
El sistema no recurre a reconocimiento facial, pero sí individualiza a las personas: les asigna un
identificador anónimo, las sigue dentro de cada cámara y las vuelve a reconocer entre cámaras
por la apariencia de su cuerpo. Esa capacidad sitúa al proyecto de lleno bajo la Ley 29733 de
Protección de Datos Personales y su reglamento vigente (Decreto Supremo 016-2024-JUS), y el
informe la documenta tal como está implementada, sin presentarla como un tratamiento
agregado que no es.

El marco teórico revisa los trabajos que sostienen el prototipo. CSRNet, P2PNet y P2R se
estudian como antecedentes del conteo de multitudes, que orientaron la discusión sobre la
granularidad de la salida, aunque el prototipo final no usa ninguno de los tres. El sistema
construido descansa en tres piezas: un detector de personas YOLO de la familia YOLO11
(Jocher y Qiu, 2024), cuya variante se elige según el hardware disponible; un seguimiento por
cámara que aplica la asociación en dos etapas de ByteTrack (Zhang et al., 2022) sobre el punto
de los pies, reforzada con señales de BoT-SORT (Aharon et al., 2022); y un descriptor de
apariencia OSNet (Zhou et al., 2019) que sostiene la reidentificación entre cámaras.

El prototipo funciona de punta a punta en un equipo sin GPU dedicada: configura el plano y
las cámaras, calibra cada cámara con una homografía, detecta y sigue a las personas, las asocia
entre cámaras, emite alertas de aglomeración con aviso por correo y produce reportes de
ocupación, flujo y oportunidad comercial. Su evaluación disponible es limitada y así se declara:
la identidad entre cámaras se midió sobre dos cámaras de demostración con cuatro personas y
veinticuatro segundos de video, con umbrales elegidos sobre esos mismos datos, de modo que
sus cifras son una demostración de funcionamiento y no una medida de precisión general.

La especificación de requerimientos adopta la numeración maestra del equipo: dieciocho casos
de uso, dieciocho requisitos funcionales y ocho no funcionales, con su matriz de trazabilidad.
El requisito no funcional rector, RNF-01, prohíbe derivar del video rasgos que permitan
identificación biométrica o reidentificación individual persistente. Con la configuración por
defecto del prototipo \textbf{ese requisito no se cumple}: la memoria de apariencia, activa salvo
que se desactive, guarda en el disco del equipo el vector OSNet de cuerpo entero de cada
persona durante veinticuatro horas (configurable hasta siete días) y la reconoce entre sesiones y
entre días. La corrección está identificada y depende solo del equipo: desactivar esa memoria
por defecto, de modo que el vector viva únicamente en memoria durante la sesión. A ese
hallazgo se suma el incumplimiento de la medida de salida agregada (M-03), porque las
grabaciones de cada sesión conservan posiciones e identificadores por persona sin un plazo de
retención automático.

Las recomendaciones plantean cerrar primero esos dos hallazgos, formalizar la Evaluación de
Impacto en Protección de Datos cuyo contenido sustantivo recoge este informe, validar la
exactitud del sistema con material del propio terminal y, solo después, considerar el despliegue
en infraestructura de nube para escalar a más cámaras simultáneas y la integración con fuentes
operativas de los locales del terminal. Ambas extensiones dependen de la autorización de LAP,
y la primera exige analizar antes el flujo transfronterizo de datos que supondría procesar fuera
del Perú (Anexo~\ref{anx:marco}).
